\begin{figure}[h]
% \setlength{\belowcaptionskip}{-20pt} 
\begin{center}

\includegraphics[width=0.99\linewidth]{Fig/Motivation.png}

\caption{Equalized odds (EO) and accuracy (Acc.) for the internal layers of the conventional ResNet on CelebA (a-b) and UTK Face (c-d) dataset. Note that lower EO represents fairer. (a) and (c) are the results using ResNet-18, while (b) and (d) use the ResNet-50. T and S stand for target and sensitive attributes, respectively. In the CelebA dataset, $a$, $e$, $m$, and $y$ represent \emph{attractiveness}, \emph{bags-under-eyes}, \emph{male}, and \emph{young}, respectively. For the UTK Face dataset, $A$ and $E$ are the \emph{Age} and \emph{Ethnicity}, respectively.
}
\label{Fig.Motivation}
\end{center}
\end{figure}

As shown in Fig. \ref{Fig.Motivation}, we observe that despite the features from a deep layer of a neural network bring high accuracy in classification, they cause fairness condition to deteriorate. We report the equalized odds (EO) and accuracy (Acc.) of ResNet-18 and ResNet-50 on the CelebA and the UTK Face dataset. We first trained a vanilla CNN, and froze the backbone. Then, we trained 3 MLP classifiers with the features from each residual module (i.e., Conv2\_x, Conv3\_x, Conv4\_x, Conv5\_x).

For CelebA dataset, in the experiments of the ResNet-18 (Fig. \ref{Fig.Motivation}(a)) and the ResNet-50 (Fig. \ref{Fig.Motivation}(b)), both the equalized odds (EO) and accuracy (Acc.) increase when the features are extracted from deeper layers. When comparing the accuracy between conv2\_x and the final layer in ResNet-18, the final layer has improved 6\% on average. However, the EO also increases over six times larger on average. As for the ResNet-50, the accuracy increased by about 16\%, and the EO increased more than eight times larger on average.

In Fig. \ref{Fig.Motivation}(c) and (d), results on UTK Face dataset again show a similar trend on the increasing EO and accuracy from shallow to deep layers. The ResNet-18 (Fig. \ref{Fig.Motivation}(c)) shows a 47\% improvement in accuracy with 1.7 times higher EO from the shallow layer (conv2\_x) to the deeper layer (Final layer). As for the ResNet-50 (Fig. \ref{Fig.Motivation}(d)), accuracy increases about 52\% and EO increases over 2.3 times.

As higher EO stands for larger bias (unfair) of the predicted result, intuitively, choosing the result at a shallow layer for final prediction could ameliorate the bias condition of the model. Our extensive experiments demonstrate that our observation could be applied to different network architectures and datasets.
